\subsection*{Additional Appendices}
\phantomsection
\addcontentsline{toc}{subsection}{Additional Appendices}

\subsubsection*{A.1\quad Wire Protocol Reference}
\phantomsection
\addcontentsline{toc}{subsubsection}{A.1\quad Wire Protocol Reference}
\label{sec:appendix-protocols}

This section provides the complete field specifications and representative implementation fragments for the three inter-component protocols described in Section~\ref{sec:wire-protocols}.

\paragraph{Complete llm-app to mcpd message reference.}

The \texttt{tool:exec} request carries the tool hash as a 64-character hexadecimal string.
If a prior \texttt{DEFER} decision returned a \texttt{ticket\_id}, the client includes it in a replay request:

\begin{verbatim}
// tool:exec request
{ "kind": "tool:exec",  "req_id": 42,
  "session_id": "sess-abc", "agent_id": "agent-1",
  "app_id": "sys_app",  "tool_id": 1,
  "tool_hash": "a3f2...c91b",   // 64-char SHA-256 hex
  "payload": {"path": "/tmp"},
  "approval_ticket_id": null }   // set on replay after DEFER

// tool:exec response (ALLOW)
{ "req_id": 42, "status": "ok",
  "result": {"os": "Linux 6.8", "cpu": "4-core"},
  "decision": "ALLOW", "reason": "admitted",
  "ticket_id": 0,  "t_ms": 3 }

// tool:exec response (DEFER)
{ "req_id": 42, "status": "error",
  "result": {}, "error": "deferred pending approval",
  "decision": "DEFER", "reason": "risk_tag:sensitive_read",
  "ticket_id": 7,  "t_ms": 1 }
\end{verbatim}

The \texttt{approval\_reply} message is issued by the client after the user confirms a deferred operation:

\begin{verbatim}
// approval_reply request
{ "sys": "approval_reply",  "session_id": "sess-abc",
  "ticket_id": 7,  "decision": "approve",
  "reason": "user confirmed",  "ttl_ms": 30000 }
\end{verbatim}

\paragraph{Complete Netlink attribute encoding.}

All attributes are 4-byte aligned.
String attributes are NUL-terminated and passed as raw bytes.
The following Python fragment shows how \texttt{mcpd} encodes a \texttt{TOOL\_REQUEST} message:

\begin{verbatim}
import struct, socket, hashlib

FAMILY_NAME = b"KERNEL_MCP\x00"
CMD_TOOL_REQUEST = 10
ATTR_AGENT_ID, ATTR_TOOL_ID = 4, 2
ATTR_REQ_ID, ATTR_TOOL_HASH = 1, 19
ATTR_AGENT_BINDING, ATTR_AGENT_EPOCH = 28, 29

def build_nla(attr_type, data):
    # NLA header: 2-byte length, 2-byte type; 4-byte aligned
    hdr = struct.pack("HH", len(data) + 4, attr_type)
    pad = (4 - len(data) % 4) % 4
    return hdr + data + b"\x00" * pad

def encode_tool_request(agent_id, tool_id, req_id,
                         tool_hash, binding, epoch):
    attrs = (
        build_nla(ATTR_AGENT_ID,      agent_id.encode() + b"\x00") +
        build_nla(ATTR_TOOL_ID,       struct.pack("=I", tool_id))  +
        build_nla(ATTR_REQ_ID,        struct.pack("=Q", req_id))   +
        build_nla(ATTR_TOOL_HASH,     tool_hash.encode() + b"\x00")+
        build_nla(ATTR_AGENT_BINDING, struct.pack("=Q", binding))  +
        build_nla(ATTR_AGENT_EPOCH,   struct.pack("=Q", epoch))
    )
    # Generic Netlink header: cmd (1), version (1), reserved (2)
    genl_hdr = struct.pack("BBH", CMD_TOOL_REQUEST, 1, 0)
    return genl_hdr + attrs
\end{verbatim}

\paragraph{Semantic hash computation.}

The semantic hash used for tool identity is a full 256-bit SHA-256 digest computed over a canonical JSON serialisation of the following manifest fields: \texttt{tool\_id}, \texttt{name}, \texttt{app\_id}, \texttt{app\_name}, \texttt{risk\_tags}, \texttt{description}, \texttt{input\_schema}, \texttt{examples}, \texttt{path\_semantics}, and \texttt{approval\_policy}.
The \texttt{endpoint} field is deliberately excluded so that redeployment to a new socket path does not invalidate the kernel registration.

\begin{verbatim}
import json, hashlib

SEMANTIC_FIELDS = [
    "tool_id", "name", "app_id", "app_name",
    "risk_tags", "description", "input_schema",
    "examples", "path_semantics", "approval_policy",
]

def semantic_hash(tool_manifest: dict) -> str:
    canonical = {f: tool_manifest[f] for f in SEMANTIC_FIELDS}
    blob = json.dumps(canonical,
                      sort_keys=True,
                      separators=(",", ":"),
                      ensure_ascii=True).encode("utf-8")
    return hashlib.sha256(blob).hexdigest()   # 64-char hex; never truncated
\end{verbatim}

The full 64-character digest is used without truncation.
A 32-bit prefix (8 hex characters) would be vulnerable to birthday collisions after approximately $2^{16}$ distinct tool registrations, which is a realistic size for a production deployment spanning multiple applications.

\paragraph{Complete manifest field specification.}

\begin{table}[htbp]
\centering
\small
\setlength{\tabcolsep}{4pt}
\begin{tabular}{l l l p{6.5cm}}
\toprule
\textbf{Field} & \textbf{Level} & \textbf{In hash?} & \textbf{Description} \\
\midrule
\texttt{app\_id}          & app  & yes & Unique application identifier string \\
\texttt{app\_name}        & app  & yes & Human-readable application name \\
\texttt{transport}        & app  & no  & Must be \texttt{"uds\_rpc"} \\
\texttt{endpoint}         & app  & no  & Absolute path under \texttt{/tmp/linux-mcp-apps/} \\
\texttt{tool\_id}         & tool & yes & Numeric identifier, unique within deployment \\
\texttt{name}             & tool & yes & Tool name presented to the LLM planner \\
\texttt{description}      & tool & yes & Natural-language description of tool behaviour \\
\texttt{risk\_tags}       & tool & yes & List of risk labels (e.g., \texttt{"sensitive\_read"}) \\
\texttt{operation}        & tool & no  & Operation string forwarded to the tool service \\
\texttt{input\_schema}    & tool & yes & JSON Schema object for payload validation \\
\texttt{examples}         & tool & yes & List of example invocation payloads \\
\texttt{path\_semantics}  & tool & yes & Path interpretation hints (e.g., \texttt{"repo\_rel"}) \\
\texttt{approval\_policy} & tool & yes & Null or policy name triggering kernel DEFER \\
\texttt{timeout\_ms}      & tool & no  & Per-invocation execution timeout \\
\bottomrule
\end{tabular}
\caption{Complete manifest field specification. Fields marked ``In hash: yes'' are included in the SHA-256 computation that produces the tool's semantic identity. Fields marked ``no'' are operational metadata excluded from the hash so that redeployment and timeout tuning do not change tool identity.}
\label{tab:manifest-fields}
\end{table}

\subsubsection*{A.2\quad System Comparison Table}
\phantomsection
\addcontentsline{toc}{subsubsection}{A.2\quad System Comparison Table}
\label{sec:appendix-comparison}

\begin{table}[htbp]
\centering
\small
\setlength{\tabcolsep}{4pt}
\begin{tabular}{p{2.1cm} p{2.4cm} p{2.4cm} p{2.4cm} p{2.2cm} p{1.8cm} p{2.4cm}}
\toprule
\textbf{System} &
\textbf{Control-plane placement} &
\textbf{Tool identity mechanism} &
\textbf{Approval / governance} &
\textbf{Failure resilience} &
\textbf{Kernel integration} &
\textbf{Security boundary} \\
\midrule

\texttt{linux-mcp} (this work)
& Kernel module (\texttt{kernel\_mcp}) + user-space gateway (\texttt{mcpd})
& Manifest-derived semantic hash, registered and verified by kernel
& Kernel-enforced approval tickets with session binding and lifecycle state
& Approval state survives \texttt{mcpd} crash (kernel holds it)
& Custom Generic Netlink module; \texttt{CAP\_NET\_ADMIN}-gated
& Kernel arbitration; enforcement independent of gateway correctness \\

\midrule

\texttt{OpenClaw} \cite{openclaw2026}
& User-space gateway (persistent process)
& Skill name and version in gateway registry
& Application-level permission checks; no durable approval state
& Lost on process restart; no persistent control-plane state
& None
& Gateway process; bypassable via logic flaws \cite{openclaw_security1, openclaw_security2} \\

\midrule

Harness Agent \cite{harness2025}
& Cloud-hosted pipeline orchestrator
& Pipeline step identity via RBAC connector
& Role-based access control, audit log, organisational policies
& Cloud-hosted; resilient by platform design
& None (SaaS)
& Platform RBAC; trust boundary is the SaaS control plane \\

\midrule

AIOS / AgentOS \cite{mei2024aios, liu2026agentos}
& New OS-level agent kernel (proposed)
& Agent and tool registration in LLM-centric scheduler
& Scheduling-layer resource governance
& Not evaluated; proposal-stage work
& New kernel substrate (proposed)
& Full OS redesign; broad scope \\

\midrule

LangChain / LangGraph \cite{langchain2022, langgraph2024}
& User-space library (in-process)
& Tool name string; no cryptographic binding
& None; developer-implemented logic
& Lost on process crash
& None
& Application process; no isolation between tool invocations \\

\midrule

Apple Intelligence \cite{apple2024pcc}
& On-device OS services + Private Cloud Compute
& System-managed capability entitlements
& OS entitlement model + PCC attestation
& High; system services are OS-persistent
& Deep OS integration (proprietary)
& Hardware attestation + OS entitlements; closed ecosystem \\

\bottomrule
\end{tabular}
\caption{Comparison of representative agent systems across control-plane and security dimensions.}
\label{tab:system_comparison}
\end{table}
